% Agent 11: Book pages 385--387 (PDF pages 22--23)
% Content: Section 2.9 Reconnection of Magnetic Field Lines (continued from p.\,384),
%          Section 3 The Origin of Stellar Black Holes,
%          Table 3.1 Summary of Numerical Computations of Supernova Explosions
%
% NOTE: The beginning of Section 2.9 and the turbulence paragraph at the
%       top of p.\,384 may overlap with Agent 10.  Only the material starting
%       from the \emph{Basic ideas} sub-heading of S2.9 onward is new here.

%% ============================================================
%%  End of Section 2.8 / transition (top of p.\,384)
%% ============================================================

uncertain state. Little is known with confidence about the astrophysical
circumstances under which turbulence should develop, or about the strength
of the turbulence in various situations. For overviews of the current state of
one's knowledge see, e.g., Chapter~3 of Landau and Lifshitz~(1959); also
Pikel'ner~(1961)\cite{Pikelner1961}.

%% ============================================================
%%  Section 2.9  Reconnection of Magnetic Field Lines
%% ============================================================

\subsection{Reconnection of Magnetic Field Lines}
\label{sec:2.9}

\textit{Basic references}\quad
\S5.3 of Cowling~(1965)\cite{Cowling1965};
Sonnerup~(1970)\cite{Sonnerup1970},
Yeh~(1970)\cite{Yeh1970},
Vainstein and Zel'dovich~(1972)\cite{VainsteinZeldovich1972}.

\bigskip\noindent
\textit{Basic ideas}

\medskip\noindent
In flat spacetime consider a plasma that is macroscopically neutral, that has a
high electrical conductivity~$\sigma$, that is sufficiently dilute for one to ignore its
dielectric properties: $\varepsilon = \mu = 1$.  Let the plasma be endowed with a large-scale
magnetic field, and examine the evolution of that field in a Lorentz frame where
the plasma has low velocity $|\mathbf{v}|\ll c$.  Maxwell's equations for the electromagnetic
field then read
%
\begin{equation}
\boldsymbol{\nabla}\times\mathbf{E} = -\boldsymbol{\nabla}\cdot\mathbf{B} = 0,
\tag{2.9.1}
\end{equation}
%
\begin{equation}
\boldsymbol{\nabla}\times\mathbf{B} - (1/c)(\partial\mathbf{B}/\partial t) = 4\pi\mathbf{J}/c.
\tag{2.9.2}
\end{equation}

In the local rest frame of the plasma the current is proportional to the electric
field, $\mathbf{J} = \sigma\mathbf{E}$.  When transformed to the Lorentz frame where the plasma moves
with velocity~$\mathbf{v}$, this equation says
%
\begin{equation}
\mathbf{A} = \sigma\!\left[\mathbf{E} + (\mathbf{v}/c)\times\mathbf{B}\right].
\tag{2.9.2}
\end{equation}

A straightforward calculation from the above equations leads to the following
law for the rate of change of magnetic field along the world lines of the plasma:
%
\begin{equation}
\frac{d\mathbf{B}}{d\tau}
  = \bigl(\omega + \sigma - \tfrac{2}{3}\theta\,\mathbf{1}\bigr)\cdot\mathbf{B}
  - \frac{1}{4\pi\sigma}
    \biggl(\frac{\partial^{2}\mathbf{B}}{\partial t^{2}} - c^{2}\,\nabla^{2}\mathbf{B}\biggr).
\tag{2.9.3}
\end{equation}

Here $\mathbf{1}$ is the unit 3-tensor; $\omega$, $\sigma$, and $\theta$ are the rotation, shear, and
expansion of the plasma [cf.\ eqs.~(2.5.17) and (2.5.18)]; and
%
\begin{equation}
d/d\tau \equiv \partial/\partial t + \mathbf{v}\cdot\boldsymbol{\nabla}.
\tag{2.9.4}
\end{equation}

Because of the very high electrical conductivity, one can ignore the
``wave-equation'' part of the evolution law~(2.9.3) almost everywhere in the
plasma; and the magnetic field evolves in a ``frozen-in'' manner:
%
\begin{equation}
d\mathbf{B}/d\tau
  = \bigl(\omega + \sigma - \tfrac{2}{3}\theta\,\mathbf{1}\bigr)\cdot\mathbf{B}.
\tag{2.9.5}
\end{equation}

%% --- page 385, right column ---

[Cf.\ Eq.~(2.5.20) and associated discussion.]  However, in regions of plasma
where the field has strong gradients, the ``wave-equation'' part must come into
play.  Consider, as the case of greatest importance, a magnetic field that is
chaotic with an ordered structure on some scale~$\ell_c$ (subscript~$c$ for
``cell size'').  At the

\begin{figure}[htbp]
\centering
% \includegraphics[width=\columnwidth]{fig_2_9_1}
\fbox{\parbox{0.9\columnwidth}{\centering [Figure 2.9.1 placeholder]\\[6pt]
The region in a plasma where magnetic field lines are reconnecting (schematic
picture).  In the top view,~(a), the reconnection region is stippled; the magnetic
field lines are drawn solidly; and the flow lines of the plasma are dotted.  The
innermost field line is in the process of reconnecting into the dashed form.  The
perspective view,~(b), shows the closed curve~$\mathscr{C}$ around which one integrates
to derive eq.~(2.9.6) for the EMF in the connecting region.}}
\caption{The region in a plasma where magnetic field lines are reconnecting.
(a)~Top view.  (b)~Perspective view showing the curve~$\mathscr{C}$.}
\label{fig:2.9.1}
\end{figure}

\noindent
interface between adjacent cells, the magnetic field must reverse sign (``neutral
point'' or ``neutral sheet''), and its gradients will be very large.  Hence, at the
interface, ``wave-equation'' behavior can come into play destroying the frozen-in
behavior of the field.  The result is a reconnection of magnetic field lines, as

%% --- page 386, left column ---

\noindent
depicted in Figure~2.9.1.  This reconnection gradually converts the two adjacent
cells into a single cell, reducing the overall strength of the magnetic field.

In situations of interest (\S\S4.4 and 5.2), the decrease of field strength due
to reconnection will be counterbalanced by an increase in strength due to
compression or shear of the plasma.

The reconnection process creates very strong electric fields.  In particular,
a simple application of Faraday's law
%
\begin{equation}
\oint_{\mathscr{C}} \mathbf{E}\cdot d\mathbf{l}
  = -\frac{d}{d\tau}\!\int \mathbf{B}\cdot d\mathbf{A},
\end{equation}
%
with the curve~$\mathscr{C}$ as shown in Figure~2.9.1b, shows that the EMF built up
in the reconnecting region is
%
\begin{equation}
\Delta\phi_e
  = (\text{rate of reconnection of magnetic flux})
  \equiv d\Psi/dt.
\tag{2.9.6}
\end{equation}

In situations of interest to us [\S5.12; Lynden-Bell~(1969)\cite{LyndenBell1969}],
these EMF's can be far greater than $10^{10}$~volts.  If the density of the plasma is
sufficiently low to permit long mean-free paths for charged particles, and if the
reconnecting region is sufficiently straight, then these EMF's can accelerate
particles to ultra-relativistic energies; and the particles can then radiate intense
synchrotron radiation.  In this manner regions of reconnecting field lines may be
sources of flares.

%% ============================================================
%%  Section 3  The Origin of Stellar Black Holes
%% ============================================================

\section{The Origin of Stellar Black Holes}
\label{sec:3}

\textit{Basic references}\quad
\S13.13 of ZN\cite{ZN1971};
Peebles~(1972)\cite{Peebles1972}.

\bigskip\noindent
\textit{Basic ideas and issues}

\medskip\noindent
How much of the mass of the Galaxy is in the form of black holes?  What is the
spectrum of black hole masses?  How many black holes are created per year by
stellar collapse in our Galaxy?  To these questions one has only the vaguest of
answers in 1972.

\textit{If\/} the spectrum of stellar masses at birth is the same elsewhere as in the solar
neighborhood, and always has been the same, then roughly half of the mass of the
Galaxy has been cycled through stars of $M > 2M_\odot$, and such stars are being
born at a rate of about 0.2~per year.  \textit{If\/} the maximum mass of a neutron star is
$2M_\odot$, and \textit{if\/} stars of $M > 2M_\odot$ lose a negligible amount of mass during their
evolutions and deaths, then all of the matter which goes into such stars eventually
winds up in black holes.  Thus, 5~per cent of the mass of the Galaxy might be in
black holes, and new black holes might be forming at a rate of ${\sim}\,0.2$~per
year---\textit{might}.  But very probably not, because the if's which go into this result
are probably not satisfied.

In particular, both observation and theory suggest that mass loss is very
significant during the late stages of stellar evolution and during the death throes
of massive stars.  Mass loss might be so important that black holes are rare
beasts, indeed!

Let us focus attention on mass loss during the death throes, as predicted by
the best numerical calculations to date.  But in doing so, let us keep in mind the
very primitive state of the modern computations---for example, their typical
neglect of the effects of rotation.

Modern computations conclude that stars whose masses exceed the
Chandrasekhar limit, $M\gtrsim 1.4\,M_\odot$, at the endpoint of their evolution should
explode as supernovae.  A supernova explosion can leave behind two remnants: an
expanding gas cloud (the outer part of the original star), and a remnant ``star''.
The remnant ``star'' will be a neutron star if its mass, $M_{\text{remnant}}$, is less
than ${\sim}\,2M_\odot$; but will be a black hole if its mass is greater than
${\sim}\,2M_\odot$.\footnote{In these calculations the initial core and final remnant
have the same mass: $M_\odot$ in one case, and $3M_\odot$ in another.
But since we have assumed that $M_{\text{star}} = M_{\text{core}}$, and since the models
give no mass loss, we must take $2.8\times M_\odot$ and $3\times 32 = M_\odot$ as the
masses of the ``true'' remnants.}

We summarize in Table~\ref{tab:3.1} the results of various calculations of the
supernova process.

%% ============================================================
%%  Table 3.1  (page 387 -- printed upside-down in the original)
%% ============================================================

\begin{table*}[htbp]
\centering
\caption{Summary of Numerical Computations of Supernova Explosions}
\label{tab:3.1}
\renewcommand{\arraystretch}{1.4}
\small
\begin{tabular}{l|c c|c c|p{4.5cm}}
\hline\hline
 & \multicolumn{2}{c|}{Initial Conditions}
 & \multicolumn{2}{c|}{Results}
 & \\[2pt]
\cline{2-5}\\[-8pt]
Authors
 & $\dfrac{M_{\text{core}}}{M_\odot}$
 & $\dfrac{M_{\text{star}}{}^\dagger}{M_\odot}$
 & $\dfrac{M_{\text{remnant}}}{M_\odot}$
 & \parbox{2.2cm}{\centering Envelope\\Kinetic Energy\\($10^{51}$~ergs)}
 & \centering Main Factor Regulating\\the Explosion\arraybackslash \\[4pt]
\hline
Fraley (1968)\cite{Fraley1968}
 & $\gtrsim 40$
 & $\gtrsim 100$
 & 0
 & 3
 & Oxygen detonation in the core at the end of quasistatic evolution \\[4pt]
\hline
Colgate and White (1966)\cite{ColgateWhite1966}
 & 10
 & 30
 & 1.8
 & 1.9
 & Absorption of neutrinos in the envelope \\[4pt]
\hline
Arnett (1966)\cite{Arnett1966}
 & 1.5
 & $\gtrsim 3.5$
 & 0.87
 & 0.1
 & \multirow{2}{4.5cm}{Carbon detonation in the core at the end of quasistatic evolution; degenerate core} \\
 & 1.4
 & 4
 & 0.1
 & 0.1
 & \\[4pt]
\hline
Ivanova, Imshennik, and
Nadezhin (1969)\cite{IvanovaImshennikNadezhin1969}
 & 2
 & $\gtrsim 4$
 & 2.46, 0.95
 & 196
 & \multirow{2}{4.5cm}{Absorption of neutrinos in the envelope} \\
 & 4
 & $\gtrsim 8$
 & 0.56, 0.66
 & $4+$
 & \\[4pt]
\hline
Nadezhin (1969)\cite{Nadezhin1969}
 & 10
 & 30
 & 9.75
 & 0.025
 & Oxygen detonation during the envelope during collapse \\[4pt]
\hline
Hansen and Barkat (1969)\cite{HansenBarkat1969}
 & 1.4
 & 4
 & $\lesssim 0$
 & 0.07
 & Carbon detonation during collapse \\[4pt]
\hline
Arnett (1969)\cite{Arnett1969}
 & 33
 & $\lesssim 100$
 & 0
 & 0
 & \parbox{4.5cm}{No mass loss; $\nu$-neutrinos. The core is opaque to absorption of neutrinos.} \\[4pt]
\hline
Arnett (1967)\cite{Arnett1967}
 & 8
 & 30
 & 9.75
 & 0
 & \parbox{4.5cm}{No mass loss; $\nu$-neutrinos. The core is opaque to absorption of neutrinos.} \\[4pt]
\hline\hline
\end{tabular}

\medskip
\begin{flushleft}
\footnotesize
$\dagger$~In recent calculations by Bruenn~(1972)\cite{Bruenn1972} the thermal
explosion of a star of $M_\odot$ produced a remnant.\\[2pt]
$\ddagger$~This does not change the conclusions given in the text about the
formation of black holes.
\end{flushleft}
\end{table*}

%% ============================================================
%%  Continuation (page 388) -- discussion of Table 3.1
%% ============================================================

In Table~3.1 $M_{\text{core}}$ is formally the mass of the entire star used in the
numerical calculations.  However, all the calculations assumed a homogeneous
initial model (typically polytropic or isothermal), whereas the theory of stellar
evolution predicts a highly inhomogeneous structure for presupernova stars.  In
particular, stellar evolution predicts a central core which is more or less
homogeneous, surrounded by a huge diffuse envelope of iron, hydrogen, and other
elements.  Thus, we must imagine the initial stars of the supernova calculations to
be the central, homogeneous core; and we must regard those calculations as
neglecting the envelope.  Such neglect is reasonable when one studies supernova
dynamics, because the core is collapsing and reexploding.  The envelope will
generally be thrown off, with little expenditure of work, by the reexploding parts
of the core.  The mass of the envelope might be twice that of the core---or
perhaps only the same as the core, or perhaps even less.  The theory of stellar
evolution is not definite on this point.  Hence the question marks in column~2 of
Table~3.1.

It is clear, from the diverse results shown in Table~3.1, that the theory of
supernova explosions is far from perfect in 1972.  Nevertheless, one can form the
following very tentative conclusions.  (i)~For stars which approach the ends of
their quasistatic evolution with masses less than ${\sim}\,12$ to $30\,M_\odot$, the
supernova explosion may produce a neutron star.  (ii)~For stars with masses
greater than ${\sim}\,12$ to $30\,M_\odot$, the explosion may produce a black hole.  If
this tentative conclusion is correct, then no more than ${\sim}\,1$~per cent of the
mass of the Galaxy should be in the form of black holes today; and new black
holes should be created at a rate no greater than ${\sim}\,0.01$~per year.
